\begin{proposition}[四阶密度累积量的完全显式式与三阶符号翻转结构]\label{prop:fold-gauge-anomaly-bernoulli-p-kappa4-and-kappa3-signflip}
\leanverified{paper\_fold\_gauge\_anomaly\_bernoulli\_p\_kappa4\_and\_kappa3\_signflip}
在定理 \ref{thm:fold-gauge-anomaly-bernoulli-p-laws} 的记号下，记 \(P_{G,p}(\theta)=\log\rho(Q_p(e^\theta))\)。
则：
\begin{enumerate}
  \item 四阶密度累积量具有完全显式闭式
  \[
  \boxed{\;
  P_{G,p}^{(4)}(0)=\frac{p^2(1-p)\,\Pi_{17}(p)}{(p+1)^7(p^2-p+1)^7}\;},
  \]
  其中
  \[
  \Pi_{17}(p):=
  \begin{aligned}[t]
  &609p^{17}-7572p^{16}+16148p^{15}-30024p^{14}+126169p^{13}-182509p^{12}+114342p^{11}\\
  &\quad-197084p^{10}+260364p^9-65334p^8-20148p^7-40612p^6+12405p^5+19432p^4\\
  &\quad-3652p^3-2466p^2+175p+81.
  \end{aligned}
  \]
  \item 多项式 \(\Pi_{11}(p)\) 在区间 \((0,1)\) 内恰有两个零点 \(p_-<p_+\)，因此 \(P_{G,p}^{(3)}(0)\) 在 \((0,1)\) 内发生两次符号翻转。
  数值上
  \[
  p_-\approx 0.3929567984354658,\qquad
  p_+\approx 0.7651253794543234.
  \]
\end{enumerate}
\end{proposition}
\begin{proof}[证明要点]
由定理 \ref{thm:fold-gauge-anomaly-bernoulli-p-laws}(4) 的 Perron 分支 \(\lambda(u,p)=\rho(Q_p(u))\) 与恒等式
\(
P_{G,p}(\theta)=\log\lambda(e^\theta,p)
\)，对四次特征方程 \(\det(\lambda I-Q_p(u))=0\) 在 \((u,\lambda)=(1,1)\) 处作四阶隐式求导并整理为 \(\theta\)-导数，即得 \(\kappa_4(p)\) 的有理闭式；
对分母作因式化可化为所示 \((p+1)^7(p^2-p+1)^7\) 的结构。
关于 \(\Pi_{11}\) 在 \((0,1)\) 内零点个数的断言，可由 Sturm 序列对根计数得到（并可据此确定 \(\kappa_3\) 的符号翻转次数）；所列数值由高精度求根给出。证毕。
\end{proof}

\endinput
